\Introduction

\textbf{Актуальность темы}. В современной практике обеспечения качества сложных программных комплексов проведение нагрузочного тестирования требует создания стабильной и управляемой среды. Одной из ключевых проблем является необходимость имитации внешних систем, с которыми взаимодействует объект тестирования. Использование имитационных сервисов («заглушек») позволяет изолировать тестируемую систему, однако управление большим количеством таких сервисов, их развертывание на удаленных серверах и контроль их состояния в ручном режиме существенно замедляют процесс подготовки испытаний.

В условиях реализации политики импортозамещения в Российской Федерации актуальность приобретает разработка отечественных инструментов управления инфраструктурой, способных заменить зарубежные программные продукты и обеспечить полноценное функционирование стендов нагрузочного тестирования в среде защищенных операционных систем, таких как ОС Astra Linux. Таким образом, создание системы, автоматизирующей жизненный цикл имитационных сервисов и предоставляющей механизмы контроля трафика, является своевременной и практически значимой задачей.

\textbf{Цель работы} -- проектирование и разработка автоматизированной информационной системы для централизованного управления распределенной инфраструктурой имитационных сервисов.

\textbf{Задачи исследования:}
\begin{compactlist}
  \item Провести анализ предметной области и существующих подходов к управлению имитационными средами в процессах обеспечения качества ПО.
  \item Спроектировать архитектуру программного комплекса, обеспечивающую гибкое масштабирование и работу в различных режимах.
  \item Обосновать выбор технологического стека и программных средств реализации системы.
  \item Разработать алгоритмы динамического проксирования трафика, механизмы контроля интенсивности запросов и программные интерфейсы для автоматизации управления.
  \item Реализовать систему хранения конфигураций и контроля состояний исполняемых файлов.
  \item Провести тестирование разработанного комплекса в среде ОС Astra Linux и оценить эффективность реализованных механизмов мониторинга и управления.
\end{compactlist}

\textbf{Объект исследования} -- инфраструктура запуска и эксплуатации имитационных сервисов (заглушек).

\textbf{Предмет исследования} -- механизмы и алгоритмы автоматизированного управления и контроля состояния имитационных сервисов.

Научная новизна работы заключается в разработке алгоритма централизованного управления распределенной средой, поддерживающего индивидуальную настройку ограничений интенсивности запросов (Rate Limiting) для каждого сервиса, а также в реализации программного интерфейса (API) для автоматического контроля состояния заглушек непосредственно из скриптов нагрузочного тестирования.

Практическая значимость работы состоит в создании программного инструмента, позволяющего отказаться от использования зарубежных систем управления, упростить эксплуатацию инфраструктуры имитационных сервисов и повысить уровень автоматизации процесса нагрузочного тестирования.

\textbf{Положения, выносимые на защиту:}
\begin{compactlist}
  \item Архитектура программного комплекса, поддерживающая гибкое масштабирование системы от локального запуска на одном хосте до формирования распределенного кластера.
  \item Механизм динамического проксирования запросов, реализующий прозрачную маршрутизацию входящего трафика к целевым имитационным сервисам.
  \item Алгоритм контроля интенсивности запросов (Rate Limiting), обеспечивающий стабильность работы имитационной среды при пиковых нагрузках.
  \item Методика формирования и экспорта метрик производительности для непрерывного мониторинга состояния инфраструктуры заглушек.
\end{compactlist}

\textbf{Структура работы.} Выпускная квалификационная работа включает введение, три главы основного текста, заключение, список использованных источников и приложения.
Во введении обоснована актуальность исследования, поставлены цель и задачи, определены объект и предмет работы, указаны научная новизна и практическая значимость полученных результатов.
В первой главе проводится анализ предметной области, исследуются существующие аналоги и методы управления имитационными средами. На основе сравнительного анализа обосновывается выбор технологий и формируются требования к системе.
Во второй главе описывается проектирование архитектуры системы, структуры хранения данных и логики работы основных модулей. Приводятся схемы взаимодействия компонентов и алгоритмы реализации функций проксирования и управления кластером.
В третьей главе рассматриваются вопросы практической реализации, настройки и эксплуатации системы в среде ОС Astra Linux. Приводятся результаты тестирования, подтверждающие корректность работы разработанных механизмов управления и мониторинга.
В заключении сформулированы основные выводы, подтверждающие решение поставленных задач и достижение цели работы.
